\section{Symmetries of the local logarithm}
\label{subsec:rgnl-symmetries}

Apply \cref{prop:rgnl-cluster-logarithm} with
$V=F$, $w=h/4$ and denote its output by
$E_k^\#(m,\mathscr E)$. Each connected term assigned to
$Y$ is a product of integrated activities on subpolymers
of $Y$. Thus ordinary and admissible pair support, local
connection dependence, and dependence only on $m'$ and
$\mathscr E^r(X)$, $X\subseteq Y$, pass to the sum.
The uniform bound and the polynomial approximants prove
analyticity of its coefficients in $(A,\mathscr C)$.

Let $\mathbb G_k^L$ denote the chamber symmetry group with
translations in $LM\Z^3$, so that it preserves every
parameter partition. Its actions and the gauge,
charge-conjugation and reality actions are those of
\cref{subsec:rg-activity-spaces}.

\begin{proposition}\label{prop:rgnl-symmetries}
The local logarithm satisfies
\begin{align}
 a_{\tilde g}E_k^\#(\tilde gY;\tilde gA,\tilde g\mathscr C)
          &=E_k^\#(Y;A,\mathscr C),\qquad \tilde g\in\mathbb G_k^L,
                                      \label{eq:rgnl-lattice-symmetry}\\
 G_{\lambda+\lambda'}E_k^\#(Y;A^1,\mathscr C^1)
          &=E_k^\#(Y;A,\mathscr C),    \label{eq:rgnl-gauge-symmetry}\\
 \mathrm{Cg}\,E_k^\#(Y;-A,-\mathscr C)
          &=E_k^\#(Y;A,\mathscr C).    \label{eq:rgnl-charge-symmetry}
\end{align}
In the gauge identity both pairs belong to the domain and
 $A^1=A-d\lambda$, $\mathscr C^1=\mathscr C-d\lambda'$
on all bonds at distance at most $1$ from $Y$, with
real $\lambda,\lambda'$. Here $d\lambda$ denotes the
lattice gradient, as in the preceding gauge conventions.
For real $m$ and $\mathscr E\in K_k^x$ one also has
\begin{equation}\label{eq:rgnl-reality-symmetry}
 \mathbf K E_k^\#(Y;-\overline A,-\overline{\mathscr C})
                                        =E_k^\#(Y;A,\mathscr C).
\end{equation}
Every monomial is balanced.
\end{proposition}

\begin{proof}
We first record how transformations pass through the
finite operations. Each of the four algebra maps is a
continuous ring automorphism, linear for the first three
and antilinear for $\mathbf K$. If an odd field transforms
like the generators it replaces, the two compositions of
the field substitution with the automorphism agree on
each generator; hence they agree on the algebra.
The same observation applies to adding a field.
Real-coefficient finite differences
\eqref{eq:rgnl-finite-difference} then commute with these
actions. A lattice symmetry also permutes the parameter
cubes and localization summands. It preserves face
connectedness, incidence relations and the Ursell coefficient.

The contour functional commutes with the linear actions.
For the antilinear action, parametrize $t=r\exp(i\theta)$:
\[
 \mathcal C_rf=\frac1{2\pi}\int_0^{2\pi}
             \frac{f(r\exp(i\theta))}{r\exp(i\theta)-1}\,\dd\theta.
\]
Conjugation of coefficients, together with
$t\mapsto\bar t$ and $\theta\mapsto-\theta$, gives the
same integral at the transformed other arguments.
The integration over $W$ commutes with every linear or
antilinear action which sends
$W\mapsto UW$, $\bar W\mapsto(U^{-1})^T\bar W$.
Indeed Wick's rule preserves every paired moment under
this change, and all these moments are real.
This applies to the lattice and gauge actions, to
$\mathbf K$ with $U=\varkappa$, and to charge conjugation
extended by the identity on $W$.
The last extension will require a separate covariance
argument below.

For a lattice symmetry the primitive covariance gives
\[
 Z(\tilde gs')[\tilde g\zeta]=\tilde gZ(s')[\zeta],\qquad
 H_k(\tilde gs'',\tilde gB)\tilde g=\tilde gH_k(s'',B),
 \qquad
 \Gamma_k(\tilde gs'',\tilde gB)\tilde g=\tilde g\Gamma_k(s'',B).
\]
The barred action is $(S^{-1})^T$ on spinors;
$(S^{-1})^T=\pm\mathsf C^{-1}S\mathsf C$ gives the
same covariance for $\bar H_k$.
H\"older quotients commute with the simultaneous
transformation of their endpoints and connection.
Thus $J_k^\#\Psi$ and $W_k^\#$ transform like their
ordinary and pair generators.
The quadratic primitive satisfies
$M_k(\tilde gX,\tilde gB)=\tilde gM_k(X,B)\tilde g^{-1}$,
so its bilinear form is invariant.
The mass form is invariant by the change of chamber
site and spin indices, and $R\mathscr E$ has the
incoming lattice symmetry. The reconstruction substitution
commutes with the action by
\eqref{eq:rg-step-substitution-covariances}.
The box measure is invariant under $\zeta\mapsto\tilde g\zeta$.
Apply these identities to $G_1,G_2,G_3$, the first
localization and substitution, then to the second
localization, $K,H$, and the convergent connected sum.
They prove \eqref{eq:rgnl-lattice-symmetry}.

For the gauge identity set $\Lambda=\lambda+\lambda'$.
In a summand assigned to $Y$, all parameters outside $Y$
vanish. The local row result proves that $B^1$ may be
replaced by $B-d\Lambda$ in every row used in that
summand. The latter field is still in the operator domain,
since real gauge shifts preserve it.
Telescoping the gradient along each path gives
\[
 U_{B-d\Lambda}(x,x')
       =U_B(x,x')\exp\big(-ie(\Lambda(c')-\Lambda(c))\big).
\]
The ordinary primitive row obeys
\[
 H_k(s'',B^1)(x,\cdot)
       =\exp(ie\Lambda(c))H_k(s'',B)(x,\cdot)
                                                 \exp(-ie\Lambda^u),
\]
when it is nonzero, where $\Lambda^u$ is the function
on the unit indices. The barred row has the inverse
phases. In the H\"older row the phase at its second
point cancels that from the transport, leaving exactly
$\exp(ie\Lambda(c))$, the phase of the first point.
The row of $\Gamma_k$ has the corresponding two-sided
gauge covariance. Its composition with $H_k^\#$ thus
gives the same covariance for $W_k^\#$, including the
rows which vanish outside $Y$.
The primitive bilinear and the local mass form are
gauge invariant, and so is $R\mathscr E$ under the
separate shifts of its two slots. The reconstruction
substitution also commutes with this local gauge action.
The photon fields and photon measure are unchanged.
The finite operations and integration therefore prove
\eqref{eq:rgnl-gauge-symmetry}.

For reality, the photon primitives are real on real
parameter families. Their analytic dependence extends
this identity to
$\overline{O(\bar s')}=O(s')$ on the full polydisk.
Since $C$ is real,
\[
 Z(\bar s')[-\zeta]=-\overline{Z(s')[\zeta]}.
\]
Transform $s\mapsto\bar s$, $t\mapsto\bar t$,
$\zeta\mapsto-\zeta$ and
$(A,\mathscr C)\mapsto(-\bar A,-\overline{\mathscr C})$.
The fermion identities are
\[
 \overline{H_k(\bar s'',-\bar B)}
                         =\varkappa^{-1}H_k(s'',B)\varkappa,\qquad
 \overline{\Gamma_k(\bar s'',-\bar B)}
                         =\varkappa^{-1}\Gamma_k(s'',B)\varkappa.
\]
For the barred field use
$\varkappa\mathsf C=-\mathsf C\varkappa$.
The transport relation
$\overline{U_{-\bar B}(x,x')}=U_B(x,x')$
then gives the same identities for the H\"older rows.
The local quadratic primitive obeys the corresponding
identity, and $\varkappa^T\varkappa=1$ preserves the
mass form. The number $m'$ is real for the present data,
and $R\mathscr E$ and the reconstruction substitution
have the reality symmetry. The contour argument just
given, invariance of the box under sign reversal, and
weighted integration prove
\eqref{eq:rgnl-reality-symmetry}.

Charge conjugation gives the first-localization symmetry
by
\[
 \begin{gathered}
 H_k(s'',-B)\mathsf C=\mathsf C\bar H_k(s'',B),\qquad
 \bar H_k(s'',-B)\mathsf C=\mathsf C H_k(s'',B),\\
 M_k(X,-B)^T=\mathsf C^{-1}M_k(X,B)\mathsf C .
 \end{gathered}
\]
Anticommuting the two transformed legs verifies the
bilinear and mass-form identities. The first-step
vertex has this symmetry by its explicit two-half formula.
It follows that
\[
 \mathrm{Cg}\,E_{\le3}(X;-A,-\mathscr C,-\zeta)
                                        =E_{\le3}(X;A,\mathscr C,\zeta).
\]
Direct covariance of the second substitution is not
available, since its unbarred and barred fluctuations
use different factors. We instead compare their Wick
covariances.

Introduce independent generators $\phi^\#$ indexed like
$\psi^\#$, with their own charge-conjugation action,
and let $\operatorname{ad}$ add them to the old generators.
Put
\begin{align*}
 \mathcal V(X)
   ={}&\operatorname{ad}E_{\le3}(X)
      +\operatorname{ad}\mathscr E^r(X)-\mathscr E^r(X)\\
     &-\mathbf1_{X=\square}m'
          \big(\operatorname{ad}
                         \langle\bar\psi,\psi\rangle_\square
                                -\langle\bar\psi,\psi\rangle_\square\big).
\end{align*}
It has charge-conjugation symmetry under
$(A,\mathscr C,\zeta)\mapsto(-A,-\mathscr C,-\zeta)$.
Let $I_{s'',B}$ fix the old generators and replace
$\phi^\#$ by $W_k^\#(s'',B)$. The second-localization
input is exactly $I_{s'',B}\mathcal V(X)$, since
the polynomial contours in $G_4,G_5$ give their value
at $1$ minus their value at $0$.
For any union $U$, interpolation restricted to cubes
of $U$, with the disconnected terms removed, gives
\begin{equation}\label{eq:rgnl-subvolume-substitution}
 \sum_{Y\subseteq U}F(Y)=
                        I_{\mathbf1_U,B}\sum_{X\subseteq U}\mathcal V(X).
\end{equation}

At $s''=\mathbf1_U$ define
\[
 \phi_\pm=H_k(\mathbf1_U,\pm B)\Gamma_k(\mathbf1_U,\pm B)W,
 \qquad
 \bar\phi_\pm=\bar H_k(\mathbf1_U,\pm B)\bar W,
 \qquad
 \phi'=\mathsf C\bar\phi_-,\quad
 \bar\phi'=\mathsf C^{-1}\phi_-.
\]
The substitution by $(\phi',\bar\phi')$ and their
H\"older quotients, denoted $I'$, satisfies
$\mathrm{Cg}\,I_{\mathbf1_U,-B}=I'\mathrm{Cg}$
on every generator, when $\mathrm{Cg}$ fixes $W$.
The only nonzero ordinary pair expectations for the
positive substitution are the kernel of
\[
 \mathcal G_U(B)
        =H_k(\mathbf1_U,B)\Gamma_k(\mathbf1_U,B)
                                            \bar H_k(\mathbf1_U,B)^T.
\]
The identity moment has kernel
$\mathbf1_{u=u'}/\nu(u)$, with $\nu=1/48$ at zero.
This precisely cancels one of the two measures in
the kernel composition, so the same formula holds
at every level. The primitive transpose relation gives
\[
 \mathcal G_U(-B)^T
                      =-\mathsf C^{-1}\mathcal G_U(B)\mathsf C^{-1}.
\]
Anticommuting $\bar\phi_-$ and $\phi_-$ therefore gives
\[
 \int\dd\mu_I(W)\,\phi'(x)\bar\phi'(x')^T
     =\mathsf C[-\mathcal G_U(-B)^T(x,x')]\mathsf C
     =\mathcal G_U(B)(x,x').
\]
Same-bar pair expectations vanish in both substitutions.
The pair generators are obtained by the identical
linear H\"older maps at $B,-B$, so their pair expectations
also agree. Wick's theorem now proves
$\int\dd\mu_I(W)I'Q=\int\dd\mu_I(W)I_{\mathbf1_U,B}Q$
for every polynomial $Q$ in the old and $\phi^\#$
generators: each monomial integral is the Pfaffian of
its pair expectations, with old generators as coefficients.

Apply this equality to the exponential of
$\sum_{X\subseteq U}\mathcal V(X)$ and use
$\zeta\mapsto-\zeta$. By
\eqref{eq:rgnl-subvolume-substitution}, the integrated
subvolume expression
$\Xi_U=\int\dd\mu_I(W)\dd\widehat\mu(\zeta)
\exp(\sum_{Y\subseteq U}F(Y))$
has charge-conjugation symmetry.
But $\Xi_U=\sum_{Z\subseteq U}H(Z)$, by the finite
Mayer expansion. Inverting this finite subset sum gives
\[
 H(Z)=\sum_{U\subseteq Z}(-1)^{|Z\setminus U|_M}\Xi_U.
\]
Thus each $H(Z)$ has the symmetry, and the real
coefficients and absolute convergence of the connected
series prove \eqref{eq:rgnl-charge-symmetry}.

Finally take a constant $\lambda$ and $\lambda'=0$
in \eqref{eq:rgnl-gauge-symmetry}. A monomial with
$n_u$ unbarred and $n_b$ barred legs acquires the factor
$\exp(ie\lambda(n_u-n_b))$. Invariance for every real
$\lambda$ and independence of the basis monomials force
its coefficient to vanish unless $n_u=n_b$.
\end{proof}

